\begin{ouverture}[college3]{15}{Puissance et énergie électrique}{Partie 3 · L’énergie, ses transferts et ses conversions}{Une lampe puissante consomme-t-elle toujours plus ?}
\objectif{Distinguer puissance et énergie.}
\objectif{Calculer avec E = P t et des unités cohérentes.}
\objectif{Comparer des usages sur une durée et un service identifiés.}
\prerequis{Proportionnalité, tension et intensité.}
\end{ouverture}
\pcPage{Le cours}{Puissance et énergie électrique}
\phantomsection\label{ess-15}
\essSection{1}{Puissance et énergie : deux grandeurs}
La puissance $P$ indique une énergie transférée par unité de durée. Son unité est le watt (W) : $1\unite W=1\unite{J/s}$. Une grande puissance n’implique pas une grande énergie totale si la durée est très courte. Il faut distinguer l’énergie disponible dans une réserve, la puissance de fonctionnement et la durée possible. Les relations sont exploitées avec des unités cohérentes et des hypothèses annoncées.\par
\pcFigure{figures/c17-puissance.pdf}{Une même énergie peut être transférée rapidement à forte puissance ou plus lentement à faible puissance.}
\essSection{2}{Relier énergie et durée}
Pour une puissance constante, $E=P\Delta t$. Avec $P$ en W et $\Delta t$ en s, $E$ est en J. Avec $P$ en W et la durée en h, l’énergie est en Wh. $1\unite{Wh}=3\,600\unite J$ et $1\unite{kWh}=1\,000\unite{Wh}$. Ne pas écrire kW/h à la place de kWh : ce ne sont pas les mêmes grandeurs.\par
\clearpage
\essSection{3}{Relier puissance, tension et intensité}
En courant continu, pour un récepteur dans les conditions du modèle, la puissance électrique reçue vaut $P=UI$, avec U en V et I en A. Les indications nominales d’un appareil décrivent un fonctionnement prévu ; elles ne doivent pas être dépassées au hasard. Une résistance peut dissiper de l’énergie thermiquement et chauffer. La valeur maximale admissible du composant compte autant que sa résistance.\par
\pcFigure{figures/c17-energie.pdf}{Le graphe E en fonction du temps est une droite pour chaque puissance constante ; sa pente dépend de P.}
\essSection{4}{Comparer des usages}
Comparer une consommation demande la puissance et la durée d’utilisation, mais aussi un service rendu comparable. Un coût peut être estimé par $\text{énergie en kWh}\times\text{prix par kWh}$, si le prix est fourni ; les tarifs des exercices sont fictifs, pas des prix actuels. La veille, la fabrication et l’alimentation du numérique peuvent ajouter des dépenses hors du simple temps d’usage affiché.\par
\begin{pcEncadre}{method}{Un exemple expliqué}Une lampe de $8\unite W$ fonctionne $3\unite h$. $E=8\times3=24\unite{Wh}=0{,}024\unite{kWh}$. En joules : $24\times3\,600=86\,400\unite J$. Avec le tarif fictif de $0{,}25$ euro par kWh, le coût d’énergie seul est $0{,}024\times0{,}25=0{,}006$ euro ; il ne représente pas le prix total d’un abonnement.\end{pcEncadre}
\pcPage{Expérimenter et retenir}{Puissance et énergie électrique}
\phantomsection\label{activite-15}
\begin{pcEncadre}{method}{Investigation en classe}\textbf{Mission documentaire.} Comparer deux éclairages assurant un service annoncé équivalent. Relever les puissances et fixer une même durée d’usage. Calculer les énergies en Wh puis kWh. Étudier un deuxième scénario où les durées diffèrent. \textbf{À rendre :} hypothèses, tableau de calcul, unités et conclusion. Si un matériel scolaire permet une mesure, le professeur choisit une alimentation protégée et un dispositif compatible ; aucun appareil domestique n’est démonté.\end{pcEncadre}
\pcSousTitre{Ce que je dois savoir expliquer}\begin{questions}
\item La puissance s’exprime en W et l’énergie en J, Wh ou kWh.
\item À puissance constante, E = P t.
\item Les unités du produit déterminent celles de l’énergie.
\item En courant continu, P = U I dans les conditions étudiées.
\item Une comparaison équitable précise durée et service rendu.
\end{questions}
\begin{pcEncadre}{attention}{Une erreur à éviter}« 2 kW pendant une demi-heure font 4 kWh. » Il faut multiplier : $2\times0{,}5=1\unite{kWh}$.\end{pcEncadre}
\pcRepere{cor-activite-15}{Pistes pour l’investigation}\pcGrille
\clearpage\markright{Je démarre}
\rubrique{Je démarre}
\par\noindent\begin{minipage}{\linewidth}
\begin{exercice}[1]{Deux grandeurs}
\begin{questions}
\item Une étiquette indique 10 W. S’agit-il d’une énergie ou d’une puissance ?
\item Quelle unité conviendrait à une énergie ?
\end{questions}
\end{exercice}
\end{minipage}
\par\noindent\begin{minipage}{\linewidth}
\begin{exercice}[1]{Une heure}
\begin{questions}
\item Une lampe de $5\unite W$ fonctionne $1\unite h$. Quelle énergie en Wh ?
\item En joules ?
\end{questions}
\end{exercice}
\end{minipage}
\par\noindent\begin{minipage}{\linewidth}
\begin{exercice}[1]{Convertir}
\begin{questions}
\item Convertir $1{,}5\unite{kWh}$ en Wh.
\item Convertir $30\unite{min}$ en h.
\end{questions}
\end{exercice}
\end{minipage}
\par\noindent\begin{minipage}{\linewidth}
\begin{exercice}[1]{Un produit électrique}
\begin{questions}
\item Un récepteur reçoit $6\unite V$ sous $0{,}20\unite A$. Calculer sa puissance.
\item Donner l’unité.
\end{questions}
\end{exercice}
\end{minipage}
\par\noindent\begin{minipage}{\linewidth}
\begin{exercice}[1]{La durée compte}
\begin{questions}
\item A fonctionne à 20 W pendant 1 h ; B à 10 W pendant 1 h. Lequel reçoit le plus d’énergie ?
\item L’égalité des durées est-elle utile ici ?
\end{questions}
\end{exercice}
\end{minipage}
\clearpage\markright{Je m’entraîne}
\rubrique{Je m’entraîne}
\par\noindent\begin{minipage}{\linewidth}
\begin{exercice}[2]{En secondes}
\begin{questions}
\item Un moteur de $3\unite W$ fonctionne $40\unite s$. Calculer E.
\item Pourquoi la réponse n’est-elle pas en Wh ?
\end{questions}
\end{exercice}
\end{minipage}
\par\noindent\begin{minipage}{\linewidth}
\begin{exercice}[2]{Une demi-heure}
\begin{questions}
\item Un appareil de $60\unite W$ fonctionne $30\unite{min}$. Calculer l’énergie en Wh et kWh.
\item Calculer en J.
\end{questions}
\end{exercice}
\end{minipage}
\par\noindent\begin{minipage}{\linewidth}
\begin{exercice}[2]{Retrouver une durée}
\begin{questions}
\item Une source transfère $180\unite J$ à puissance constante de $6\unite W$. Quelle durée ?
\item Vérifier.
\end{questions}
\end{exercice}
\end{minipage}
\par\noindent\begin{minipage}{\linewidth}
\begin{exercice}[2]{Une intensité nominale}
\begin{questions}
\item Un moteur scolaire de $2{,}4\unite W$ fonctionne sous $6\unite V$ en courant continu. Calculer I dans ce modèle.
\item Exprimer en mA.
\end{questions}
\end{exercice}
\end{minipage}
\par\noindent\begin{minipage}{\linewidth}
\begin{exercice}[2]{Un coût fictif}
\begin{questions}
\item Une énergie de $2{,}5\unite{kWh}$ est facturée au prix fictif de $0{,}20$ euro/kWh. Quel coût d’énergie ?
\item Cela inclut-il forcément l’abonnement ?
\end{questions}
\end{exercice}
\end{minipage}
\clearpage\markright{Je raisonne}
\rubrique{Je raisonne}
\par\noindent\begin{minipage}{\linewidth}
\begin{exercice}[3]{Puissante mais brève}
\begin{questions}
\item A reçoit $100\unite W$ pendant $1\unite{min}$ ; B $10\unite W$ pendant $20\unite{min}$. Comparer les énergies.
\item Pourquoi la puissance seule ne suffit-elle pas ?
\end{questions}
\end{exercice}
\end{minipage}
\par\noindent\begin{minipage}{\linewidth}
\begin{exercice}[3]{Une veille}
\begin{questions}
\item Un dispositif consomme $2\unite W$ en veille pendant $20\unite h$ par jour. Calculer l’énergie quotidienne puis pour 30 jours.
\item Que manque-t-il pour l’énergie totale de l’appareil ?
\end{questions}
\end{exercice}
\end{minipage}
\par\noindent\begin{minipage}{\linewidth}
\begin{exercice}[3]{Une résistance qui chauffe}
\begin{questions}
\item Une résistance de $100\unite\Omega$ est sous $5\unite V$. Trouver I puis P.
\item Une puissance admissible de 0,125 W suffirait-elle ?
\end{questions}
\end{exercice}
\end{minipage}
\par\noindent\begin{minipage}{\linewidth}
\begin{exercice}[3]{Même service}
\begin{questions}
\item Pour un éclairage annoncé équivalent, A vaut 40 W et B 8 W. Calculer l’économie sur 100 h en kWh.
\item Quelle limite ajouter à la conclusion ?
\end{questions}
\end{exercice}
\end{minipage}
\par\noindent\begin{minipage}{\linewidth}
\begin{exercice}[3]{Une unité fausse}
\begin{questions}
\item Une copie écrit « E = 2 kW/h après deux heures ». Identifier le problème.
\item Donner le résultat pour une puissance de 2 kW pendant 2 h.
\end{questions}
\end{exercice}
\end{minipage}
\pcPage{Défi du chapitre}{Puissance et énergie électrique}
\phantomsection\label{defi-ex-15}
\pcDefiTitre{15}{Le projet d’éclairage}
\begin{questions}
\item Dix lampes de $12\unite W$ fonctionnent $4\unite h$ par jour pendant 20 jours. Calculer l’énergie totale.
\item On les remplace par des lampes de 8 W assurant le même service. Calculer l’économie.
\item À un prix fictif de 0,25 euro/kWh, quel gain sur le seul poste énergie ?
\end{questions}
\pcGrille
